\section{Conclusion}\label{sec:conclusion}

We modelled self-propagating npm worms as a two-channel SIR epidemic and calibrated it to the two
Shai-Hulud waves. We validated its structure against the real victims, who cluster by shared maintainer
rather than by dependency links. We then used the model to evaluate where outbreaks start, which
ecosystem regimes change the epidemic threshold, and how to spend a limited response budget.

Dependency topology alone confines outbreaks to the dependents of a few hubs, and dependency-side
hardening raises that bound further. Stolen publish tokens remove it: with Shai-Hulud-like dynamics, a
worm that enters the high-impact core from an obscure leaf package can reach a large part of it, and
reactive patching after a day's delay barely helps. Breaking the credential channel, and proactively
protecting the most-depended-upon packages, are the most effective measures against the next worm.
